% authority: science_draft
\documentclass[11pt]{article}

\usepackage[margin=1in]{geometry}
\usepackage{amsmath,amssymb,amsthm}
\usepackage{booktabs}
\usepackage{longtable}
\usepackage{enumitem}

\newcommand{\EqSrc}{\mathrm{EqSrc}}
\newcommand{\EqSrcT}{\mathrm{EqSrc}_{T}}
\newcommand{\RetainH}{\mathrm{RetainH}}
\newcommand{\GenH}{\mathrm{GenH}}
\newcommand{\Fsrc}{\mathcal{F}_{\mathrm{src}}}
\newcommand{\Osrc}{\mathcal{O}_{\mathrm{src}}}
\newcommand{\Msrc}{\mathcal{M}_{\mathrm{src}}}
\newcommand{\Isrc}{\mathcal{I}_{\mathrm{src}}}
\newcommand{\Asrc}{\mathcal{A}_{\mathrm{src}}}
\newcommand{\Psrc}{\mathcal{P}_{\mathrm{src}}}
\newcommand{\Nsrc}{\mathcal{N}_{\mathrm{src}}}
\newcommand{\Csrc}{C_{\mathrm{src}}}
\newcommand{\Block}{\mathrm{Block}}
\newcommand{\NoTarget}{\mathsf{NoTarget}}
\newcommand{\Hret}{\rho_{H}}
\newcommand{\Hgen}{\gamma_{H}}

\newtheorem{definition}{Definition}
\newtheorem{lemma}{Boundary Lemma}
\newtheorem{witness}{Minimal Witness Slot}
\newtheorem{remark}{Boundary Remark}

\title{RetainH and GenH Primitive-Boundary Extraction v1}
\author{Ontology Formalizer Packet RT-20260707-021}
\date{2026-07-07}

\begin{document}
\maketitle

\section{Control Status}

This artifact implements v18 P3-T03 as a draft/control science artifact. It
extracts the exact primitive-boundary consequences for \(\RetainH\) and
\(\GenH\) from the P3-T02 conditional \(\EqSrcT\) family-closure theorem
candidate.

The primary result is:
\[
  \texttt{primitive\_boundary\_extracted\_no\_adoption}.
\]

\begin{verbatim}
primitive_boundary_result: primitive_boundary_extracted_no_adoption
RetainH_adopted: false
GenH_adopted: false
adoption_requested: false
next_route: P3-T04
\end{verbatim}

The artifact does not discharge general \(\EqSrc\), adopt \(\RetainH\), adopt
\(\GenH\), adopt a source law, edit canonical ontology, adopt
\(\mathrm{MetricData}(E)\), expand \(g_{\mathrm{eff}}\), derive matter
coupling, derive Einstein equations, promote a benchmark, issue a Gate Chair
verdict, authorize external outreach, or claim completed derivation.

\section{Inputs from P3-T02}

P3-T02 supplied a conditional theorem candidate over an already-declared source
family
\[
  \Fsrc=(\Osrc,\Msrc,\Isrc,\Csrc,\Asrc,\Psrc,\Nsrc).
\]
The theorem branch assumed:
\begin{enumerate}[label=(H\arabic*)]
  \item source-only objects in \(\Osrc\);
  \item identity witnesses in \(\Msrc\);
  \item inverse witnesses in \(\Msrc\);
  \item composition closure in \(\Msrc\);
  \item source invariant, annotation, and forbidden-channel ledger stability;
  \item source-only comparison by \(\Csrc\);
  \item a no-target-import guard for all witnesses and comparison premises.
\end{enumerate}

Under those supplied hypotheses, P3-T02 stated that \(\EqSrcT\) is reflexive,
symmetric, and transitive on \(\Osrc\). The theorem is conditional: it does not
prove that current ontology derives H1--H7 for every future source family.

P3-T02 also recorded the missing-inverse countermodel slot. That slot shows
that theorem hypotheses cannot be silently removed, but it does not by itself
adopt \(\RetainH\) or \(\GenH\).

\section{\(\RetainH\) Boundary}

\begin{verbatim}
section_marker: RetainH Boundary
\end{verbatim}

\begin{definition}[RetainH boundary pressure]
An \(H\)-retention extension introduces a source-side operation or relation
\[
  \Hret:X\rightsquigarrow X_{H}
\]
and asks whether source records, certificate status, invariant ledger entries,
annotation entries, forbidden-channel or proxy-edge incidence, comparison
status, and negative controls are retained from \(X\) to \(X_{H}\).
\end{definition}

The closed-family P3-T02 theorem does not require this primitive. It assumes
the relevant family, witnesses, comparison rule, and ledger preservation
hypotheses directly. Therefore \(\RetainH\) is avoided in that theorem branch.

For an \(H\)-retention extension, however, the current source record does not
derive a rule saying that \(\Hret\) preserves the entries needed by
\(\EqSrcT\). A candidate definition is therefore needed before a future packet
can claim that retention under \(H\) preserves the source equivalence evidence.

\begin{verbatim}
RetainH_status_for_closed_declared_family: not_required_here
RetainH_status_for_H_retention_extension: candidate_definition_needed
RetainH_adopted: false
\end{verbatim}

\section{\(\GenH\) Boundary}

\begin{verbatim}
section_marker: GenH Boundary
\end{verbatim}

\begin{definition}[GenH boundary pressure]
An \(H\)-generated family extension introduces a source-side generation rule
\[
  \Hgen:S\rightsquigarrow \Fsrc^{H}
\]
and asks which generated objects, morphisms, ledgers, comparison entries, and
negative controls are admitted into the source family used by \(\EqSrcT\).
\end{definition}

The closed-family P3-T02 theorem does not require this primitive. It assumes a
declared family \(\Fsrc\) and tests identity, inverse, composition, ledger,
comparison, and no-target conditions inside that family. It does not claim to
generate the family.

For an \(H\)-generated extension, the current source record does not derive a
rule saying which \(H\)-generated objects and morphisms are members of
\(\Osrc\) and \(\Msrc\), nor which generated invariant and negative-control
entries accompany them. A candidate definition is therefore needed before a
future packet can quantify over an \(H\)-generated source family.

\begin{verbatim}
GenH_status_for_closed_declared_family: not_required_here
GenH_status_for_H_generated_extension: candidate_definition_needed
GenH_adopted: false
\end{verbatim}

\section{Primitive Needed, Primitive Avoided, or Primitive Deferred Classification}

\begin{longtable}{p{0.21\linewidth}p{0.30\linewidth}p{0.35\linewidth}}
\toprule
Primitive & Classification & Reason \\
\midrule
\(\RetainH\) in the already-declared closed-family theorem & \texttt{not\_required\_here} & The theorem uses supplied identity, inverse, composition, ledger, comparison, and no-target hypotheses inside \(\Fsrc\). \\
\(\RetainH\) for \(H\)-retention extensions & \texttt{candidate\_definition\_needed} & A source-side retention rule is needed to preserve record and ledger evidence under \(\Hret\). \\
\(\GenH\) in the already-declared closed-family theorem & \texttt{not\_required\_here} & The theorem does not generate family members; it works only on the declared \(\Fsrc\). \\
\(\GenH\) for \(H\)-generated family extensions & \texttt{candidate\_definition\_needed} & A source-side generation and membership rule is needed before generated objects and morphisms can be used in \(\EqSrcT\). \\
\bottomrule
\end{longtable}

No classification in this packet is \texttt{required\_but\_missing} for the
already-declared closed-family theorem. No classification in this packet is
\texttt{countermodel\_blocks\_current\_route}. The route remains open to
P3-T04 audit and P3-T05 stress.

\section{Minimal Witness or Countermodel}

\begin{verbatim}
witness_slots:
  - apply_H_retention_without_RetainH
  - expand_source_family_without_GenH
\end{verbatim}

\begin{witness}[Applying H-retention without RetainH]
Let \(X\in\Osrc\) have source ledger entries sufficient for comparison. Let an
extension ask to compare \(X\) with \(X_{H}\), where \(X_{H}\) is produced by a
purported \(H\)-retention operation \(\Hret:X\rightsquigarrow X_{H}\). If no
source-side \(\RetainH\) rule states that \(\Hret\) preserves
\(\Isrc,\Asrc,\Psrc,\Nsrc\) and comparison status, then the positive
\(\EqSrcT\) branch is unavailable. The comparison must block or remain
undefined for the extension.
\end{witness}

This witness does not affect the P3-T02 theorem over an already-declared
closed family. It only shows why a future \(H\)-retention theorem needs a
candidate source definition before using retention as evidence.

\begin{witness}[Expanding the source family without GenH]
Let \(S=\{X\}\) be a seed source fragment. Suppose a later packet asserts a
generated family \(\Fsrc^{H}\) containing \(X\), \(Y=\Hgen(X)\), and a morphism
\(f:X\to Y\). If no source-side \(\GenH\) rule places \(Y\) in \(\Osrc\),
places \(f\) in \(\Msrc\), and supplies generated ledger and negative-control
entries, then \(Y\) and \(f\) cannot serve as \(\EqSrcT\) witnesses. The
comparison must block or remain outside the declared family.
\end{witness}

This witness does not prove that an \(H\)-generated extension is impossible. It
shows that the current closed-family theorem cannot be reused as a generated
family theorem without a separate source-side generation primitive or an
equivalent explicit construction.

\section{No-Target Import Guard}

The primitive-boundary classifications above are source-side classifications.
They cannot be repaired or strengthened by importing:
\begin{itemize}[leftmargin=*]
  \item target topology, target smooth atlas, or target metric;
  \item proper time, detector protocol, empirical calibration, or benchmark behavior;
  \item stress-energy semantics, stress-energy tensor, matter action, matter coupling, or Einstein equations;
  \item Gate Chair status, validation success, registry metadata, role identity, handoff state, commit state, generated derivative, or external reputation.
\end{itemize}

Those objects may appear as forbidden conclusions or process evidence. They
are not mathematical premises for \(\RetainH\), \(\GenH\), or \(\EqSrcT\).

\section{Distance-to-GR Effect}

\begin{longtable}{p{0.30\linewidth}p{0.56\linewidth}}
\toprule
Burden & Status after this packet \\
\midrule
Source ontology primitives & Unchanged; no canonical ontology edit. \\
Source equivalence EqSrc & Primitive boundaries extracted; no general EqSrc discharge. \\
RetainH & Not adopted; not required for the supplied closed-family theorem; candidate definition needed for H-retention extensions. \\
GenH & Not adopted; not required for the supplied closed-family theorem; candidate definition needed for H-generated family extensions. \\
ObsLoc\_lc & Unchanged. \\
Resp\_lc & Unchanged. \\
M\_src & Unchanged. \\
g\_eff & Unchanged; no MetricData(E) or metric-scope expansion. \\
matter coupling & Blocked and unchanged; no coupling derivation or adoption. \\
Einstein equations & Blocked and unchanged; no field-equation derivation. \\
finite-variation robustness & Unchanged; no global robustness theorem. \\
benchmark promotion & Blocked; no benchmark evidence or promotion. \\
Gate Chair status & Unchanged; no Gate Chair verdict. \\
current route freeze or hard-fail status & Not frozen; route continues to P3-T04 audit. \\
\bottomrule
\end{longtable}

\begin{verbatim}
distance_to_gr_delta:
  changed: false
  effect: no_distance_delta
  ledger_row_updated: false
  reason: P3-T03 extracts primitive boundaries only. It does not adopt RetainH
    GenH or EqSrc and does not prove the closure hypotheses from current
    ontology.
\end{verbatim}

\section{Forbidden Conclusions}

This artifact does not authorize:

\begin{itemize}[leftmargin=*]
  \item canonical ontology edit;
  \item general \(\EqSrc\) discharge;
  \item \(\RetainH\) adoption;
  \item \(\GenH\) adoption;
  \item source-law adoption;
  \item source detector/readout semantics adoption;
  \item \(\mathrm{MetricData}(E)\) adoption;
  \item \(g_{\mathrm{eff}}\) adoption or scope expansion;
  \item matter-coupling derivation or adoption;
  \item stress-energy semantics, stress-energy tensor, or matter action;
  \item Einstein-equation derivation;
  \item benchmark promotion;
  \item Gate Chair verdict;
  \item external outreach;
  \item completed derivation;
  \item future source-extension impossibility;
  \item global no-go conclusion.
\end{itemize}

\section{Source Materials}

The AEther-Flow Research Project. (2026a). \emph{Recommendations
implementation plan continue task v18} [Internal implementation plan].
\texttt{implementations\_plans/recommendations\_implementation\_plan\_continue\_task-v18.md}.

The AEther-Flow Research Project. (2026b). \emph{EqSrc family-closure theorem
or countermodel v1} [Research-control TeX artifact].
\texttt{research\_control/tasks/RT-20260707-020/artifacts/eqsrc\_family\_closure\_theorem\_or\_countermodel\_v1.tex}.

The AEther-Flow Research Project. (2026c). \emph{V18 P2-T05
source-equivalence typed object Refuter stress v1} [Research-control TeX
artifact].
\texttt{research\_control/tasks/RT-20260707-017/artifacts/source\_equivalence\_typed\_object\_refuter\_stress\_v1.tex}.

\end{document}
